%! Polynomial Functions and Real Zeros — Unit 2, Lesson 2.3 — Exit Ticket
\documentclass[10pt]{article}
\usepackage{precalculus-article}
\usepackage{precalculus-boxes}
\newcommand{\TallMath}[1]{$\displaystyle #1\rule[-1.4em]{0pt}{3.2em}$}

\begin{document}

\pageheader{Unit 2, Lesson 2.3}{Exit Ticket}

\vspace{0.12in}

\begin{enumerate}[itemsep=12pt]

\item Let $p(x) = (x - 4)(x + 1)^2$. Complete the table.

\smallskip
\renewcommand{\arraystretch}{1.5}
\begin{tabular}{|c|c|c|}
\hline
\textbf{Zero} & \textbf{Multiplicity} & \textbf{Crosses or tangent?} \\
\hline
$x = $ \blank{1.2cm} & \blank{1.2cm} & \blank{2.4cm} \\
\hline
$x = $ \blank{1.2cm} & \blank{1.2cm} & \blank{2.4cm} \\
\hline
\end{tabular}
\smallskip

Degree of $p$: \blank{1.6cm}

\item Write a polynomial in factored form with a zero at $x = -3$ of
      multiplicity 1 and a zero at $x = 2$ of multiplicity 2. Use leading
      coefficient 1.

      \vspace{0.05in}
      $p(x) = $ \blank{5.5cm} \qquad Its degree is \blank{1.4cm}

\item The table shows a polynomial function $g$ at equally spaced inputs.

\smallskip
\begin{tabular}{|c|c|c|c|c|c|}
\hline
$x$ & 0 & 1 & 2 & 3 & 4 \\
\hline
$g(x)$ & 1 & 4 & 11 & 22 & 37 \\
\hline
\end{tabular}
\smallskip

\begin{enumerate}[label=(\alph*), itemsep=4pt]
  \item 1st differences: \blank{1cm} , \blank{1cm} , \blank{1cm} , \blank{1cm}
  \item 2nd differences: \blank{1cm} , \blank{1cm} , \blank{1cm}
  \item The degree of $g$ is \blank{1.4cm}
\end{enumerate}

\end{enumerate}

\end{document}
